% ch7_c 第7章 §7.3 (书页 301–313 / p316–p328)
% 块界说明：
% (1) 起点——书页 300（p315）以完整段落收束（"…where p1 = m is an odd integer.
%     Equation (7.3.13) is the effective theory of a second-level hierarchical FQH state."），
%     书页 301 顶部"有效理论可用于确定……"为新段落，故本块无 %--B-TAIL-- 区、正文不以 %JOIN% 起头。
% (2) 终点——书页 313（p328）末句止于 "These charged excitations carry integer charges and"，
%     其续文为 p329（书页 314）顶部 "are created by the electron operators Ψ†.
%     The above theory of edge excitations is formulated in terms of …"。
%     按约定该半句译文（"……携带整数电荷，并且"）留在本文件正文最末，未写入任何 TAIL 区；
%     ch7_d 请以 %JOIN% 起头续译"由电子算符 $\Psi^{\dagger}$ 所产生。……"，勿重复翻译本块已有内容。
% (3) 本块覆盖：7.3.3 后半（自书页 301 顶部起）、7.3.4（含表 7.1、7.2，问题 7.3.3–7.3.6）、
%     7.4 章节首与要点、7.4.1（图 7.15、7.16，脚注 47，问题 7.4.1–7.4.3）、7.4.2 至书页 313 末。
有效理论可以用来确定分级（hierarchical）FQH 态的物理性质。总填充因子由运动方程
$\delta\mathcal{L}/\delta a_{0}\allowbreak=\delta\mathcal{L}/\delta\tilde{a}_{0}\allowbreak=0$ 确定，其中
\begin{equation*}
-eB=p_{1}b-\tilde{b},\qquad b=p_{2}\tilde{b}
\end{equation*}
我们求得
\begin{equation*}
\nu=\frac{b}{-eB}=\frac{1}{p_{1}-\frac{1}{p_{2}}}. \tag{7.3.14}
\end{equation*}
通过引入 $(a_{1\mu},a_{2\mu})=(a_{\mu},\tilde{a}_{\mu})$，式 (7.3.13) 可以写成下面更紧凑的形式：
\begin{equation*}
\mathcal{L}=-\frac{1}{4\pi}K_{IJ}a_{I\mu}\partial_{\nu}a_{J\lambda}\varepsilon^{\mu\nu\lambda}
+\frac{e}{2\pi}q_{I}A_{\mu}\partial_{\nu}a_{I\lambda}\varepsilon^{\mu\nu\lambda} \tag{7.3.15}
\end{equation*}
其中 $K$ 是整数矩阵
\begin{equation*}
K=\begin{pmatrix} p_{1} & -1\\ -1 & p_{2} \end{pmatrix}
\end{equation*}
而 $\boldsymbol{q}$ 是整数矢量 $\boldsymbol{q}^{\top}=(q_{1},q_{2})=(1,0)$。这里 $\boldsymbol{q}$ 将被称为电荷矢量（charge vector）。填充因子 (7.3.14) 可以改写为 $\nu=\boldsymbol{q}^{\top}K^{-1}\boldsymbol{q}$。当 $(p_{1},p_{2})=(3,2)$ 时，该分级态对应于实验中观察到的 $\nu=2/5$ FQH 态。

第二级分级 FQH 态包含两种准粒子。一种是原来电子凝聚体中的准空穴（或涡旋），另一种是新玻色凝聚体中的准空穴（或涡旋）。这两种准空穴分别通过插入源项 $-j^{\mu}a_{\mu}$ 与 $-\tilde{j}^{\mu}\tilde{a}_{\mu}$ 来产生，其中 $\tilde{j}^{\mu}$ 与 $j^{\mu}$ 具有和式 (7.3.2) 类似的形式。第一种准空穴通过在电子波函数上乘以 $\prod_{i}(\xi-z_{i})$ 产生，而第二种准空穴则通过在玻色 Laughlin 波函数上乘以 $\prod_{i}(\eta-\xi_{i})$ 产生（这里 $\xi_{i}$ 是玻色子的复坐标，$\eta$ 是准空穴的位置）。

一个一般的准粒子由 $l_{1}$ 个第一种准粒子和 $l_{2}$ 个第二种准粒子组成，并由两个整数标记。这样的准粒子携带 $l_{1}$ 个单位的 $a_{1\mu}$ 荷和 $l_{2}$ 个单位的 $a_{2\mu}$ 荷，由下式描述
\begin{equation*}
(l_{1}a_{1\mu}+l_{2}a_{2\mu})j^{\mu} \tag{7.3.16}
\end{equation*}
积掉规范场之后，我们发现这样的准粒子携带 $\sum_{J}K_{IJ}l_{J}$ 个单位的 $a_{I\mu}$ 通量。因此，这样的准粒子的统计由下式给出
\begin{equation*}
\theta=\pi\pmb{l}^{\top}K^{-1}\pmb{l}
=\frac{1}{p_{2}p_{1}-1}\left(p_{2}l_{1}^{2}+p_{1}l_{2}^{2}+2l_{1}l_{2}\right)
\end{equation*}
而准粒子的电荷为
\begin{equation*}
Q_{q}=-e\boldsymbol{q}^{\top}K^{-1}\pmb{l}
=-e\frac{p_{2}l_{1}+l_{2}}{p_{2}p_{1}-1}
\end{equation*}
对于 $\nu=2/5$ 态（即 $(p_{1},p_{2})=(3,2)$），具有最小电荷的准粒子由 $(l_{1},l_{2})=(0,1)$ 标记。最小电荷为 $-e/5$。这样的准粒子具有统计 $\theta=\frac{3}{5}\pi$。

我们也可以构造更一般的 FQH 态。这些 FQH 态的有效理论仍具有式 (7.3.15) 给出的形式，只是现在指标 $I$ 从 1 取到某个整数 $n$（$n$ 将被称为该 FQH 态的级（level））。为了得到矩阵 $K$ 的形式，让我们假设在第 $n-1$ 级，有效理论由 $a_{I\mu}$（$I=1,\ldots,n-1$）以及 $K=K^{(n-1)}$ 的式 (7.3.15) 给出。准粒子携带 $a_{I\mu}$ 规范场的整数荷。现在考虑第 $n$ 级分级态，它通过具有 $a_{I\mu}$ 荷 $l_{I}|_{I=1,\ldots,n-1}$ 的准粒子的"凝聚"而得到。这个第 $n$ 级分级态的有效理论将由含 $n$ 个规范场的式 (7.3.15) 给出。第 $n$ 个规范场 $a_{n\mu}$ 来自新的凝聚体。矩阵 $K$ 由下式给出
\begin{equation*}
K^{(n)}=\begin{pmatrix} K^{(n-1)} & -\pmb{l}\\ -\pmb{l}^{\top} & p_{n} \end{pmatrix}
\end{equation*}
其中 $p_{n}$ 为偶数。电荷矢量 $\boldsymbol{q}$ 仍由 $(1,0,0,\ldots)$ 给出。通过迭代我们看到，广义分级态总是由整数对称矩阵描述，其中 $K_{II}$ 为偶数，唯独 $K_{11}$ 为奇数。新的凝聚体给出一种新的准粒子，它同样携带新规范场 $a_{n\mu}$ 的整数荷。因此，一个一般的准粒子总是携带 $a_{I\mu}$ 场的整数荷。

让我们用更一般的语言总结上述结果。我们知道，一个分级（或广义分级）FQH 态包含许多不同的凝聚体。这些凝聚体并不相互独立：一个凝聚体中的粒子对其他凝聚体中的粒子而言就像一根通量管（flux tube）。为了描述这种耦合，方便的做法是用 $U(1)$ 规范场来描述第 $I$ 个凝聚体的密度和流 $J_{I\mu}$：
\begin{equation*}
J_{I}^{\mu}=\frac{1}{2\pi}\varepsilon^{\mu\alpha\beta}\partial_{\alpha}a_{I\beta} \tag{7.3.17}
\end{equation*}
这时，不同凝聚体之间的耦合由规范场的 Chern--Simons 项描述：
\begin{equation*}
\mathcal{L}=-\frac{1}{4\pi}K_{IJ}a_{I\mu}\partial_{\nu}a_{J\lambda}\varepsilon^{\mu\nu\lambda}
+\frac{e}{2\pi}q_{I}A_{\mu}\partial_{\nu}a_{I\lambda}\varepsilon^{\mu\nu\lambda} \tag{7.3.18}
\end{equation*}
当 $K$ 取为一般的整数 $n\times n$ 矩阵，且 $K_{II}|_{I=1}$ 为奇数、$K_{II}|_{I>1}$ 为偶数、$\boldsymbol{q}^{\top}=(1,0,\ldots,0)$ 时，式 (7.3.18) 描述的是最一般的（阿贝尔）FQH 态 (Wen and Zee, 1992a)。填充因子由下式给出
\begin{equation*}
\nu=\boldsymbol{q}^{\top}K^{-1}\boldsymbol{q} \tag{7.3.19}
\end{equation*}
准粒子激发可以看作不同凝聚体中的涡旋。一个一般的准粒子由 $n$ 个整数 $l_{I}|_{I=1,..,n}$ 标记，并可由下面的源项产生：
\begin{equation*}
\mathcal{L}=l_{I}a_{I\mu}j^{\mu} \tag{7.3.20}
\end{equation*}
这样的准粒子将记作 $\psi_{\pmb{l}}$。式 (7.3.20) 中的 $j^{\mu}$ 具有如下形式
\begin{gather*}
\pmb{j}(x)=\dot{\pmb{x}}_{0}\delta(\pmb{x}-\pmb{x}_{0})\\
j^{0}(x)=\delta(\pmb{x}-\pmb{x}_{0}) \tag{7.3.21}
\end{gather*}
它在 $\pmb{x}_{0}$ 处产生一个准粒子。

当我们产生一个准粒子 $\psi_{\pmb{l}}$ 时，它会在所有凝聚体的密度中引起一个改变 $\delta J_{I}^{0}$。由式 (7.3.18) 与式 (7.3.20) 的运动方程，我们发现 $\delta J_{I}^{0}$ 满足
\begin{equation*}
\int\mathrm{d}^{2}x\,\delta J_{I}^{0}
=l_{J}(K^{-1})_{JI}\int\mathrm{d}^{2}x\,j^{0}=(\pmb{l}^{\top}K^{-1})_{I} \tag{7.3.22}
\end{equation*}
准粒子 $\psi_{\pmb{l}}$ 的电荷与统计由下式给出
\begin{equation*}
\theta_{\pmb{l}}=\pi\pmb{l}^{\top}K^{-1}\pmb{l},\qquad
Q_{\pmb{l}}=-e q_{I}\int\mathrm{d}^{2}x\,\delta J_{I}^{0}=-e\pmb{l}^{\top}K^{-1}\boldsymbol{q} \tag{7.3.23}
\end{equation*}
准粒子统计的这一结果来自如下观察：$2\pi\delta J_{I}^{0}$ 是 $a_{I\mu}$ 的通量，而准粒子携带 $l_{I}$ 个单位的 $a_{I\mu}$ 荷。

一个一般的电子激发可以看作一种特殊的（一般）准粒子，$\psi_{e}=\psi_{\pmb{l}_{e}}$，其中整数矢量 $\pmb{l}_{e}$ 由下式给出
\begin{equation*}
l_{eI}=K_{IJ}L_{J},\qquad L_{I}=\text{整数},\qquad q_{I}L_{I}=1 \tag{7.3.24}
\end{equation*}
我们可以证明，这些电子激发满足下列性质：它们携带单位电荷（见式 (7.3.23)）；它们具有费米统计；让一个电子激发 $\psi_{e}=\psi_{L}$ 绕任意准粒子激发 $\psi_{\pmb{l}}$ 一周，总是诱导 $2\pi$ 整数倍的相位；并且由式 (7.3.24) 定义的激发，就是满足上述三条性质的全部激发。

由式 (7.3.15) 的一般有效理论，我们发现广义分级态可以由一个整数值的 $K$ 矩阵和一个电荷矢量 $\boldsymbol{q}$ 来标记。现在我们想问这样一个问题：不同的 $(K,\boldsymbol{q})$ 是否描述不同的 FQH 态？注意到，通过对规范场 $a_{I\mu}$ 的重新定义，总可以把 $K$ 矩阵对角化成对角元为 $\pm1$ 的矩阵。这样看来，所有具有相同符号差（signature）的 $K$ 矩阵似乎都描述相同的 FQH 态，因为经过规范场的适当重新定义后，它们给出同一个有效理论。这个结论当然是错误的。我们想强调，仅凭有效拉格朗日量 (7.3.15) 并不能恰当描述分级态的内部序（或拓扑序）。正是有效拉格朗日量 (7.3.15) 连同 $a_{I\mu}$ 荷的量子化条件一起，才刻画了拓扑序。配备了允许 $U(1)$ 荷之量子化条件的 $U(1)$ 规范理论被称为紧致 $U(1)$ 理论。我们的有效理论 (7.3.15) 实际上是一个所有 $U(1)$ 荷都量子化为整数的紧致 $U(1)$ 理论。因此，允许的 $U(1)$ 荷构成一个 $n$ 维立方格子，我们将称之为电荷格子（charge lattice）。所以，当考虑两个不同 $K$ 矩阵的等价性时，只能使用保持电荷量子化条件不变（即保持电荷格子不变）的场重新定义。把电荷格子映到其自身的变换属于群 $SL(n,Z)$（行列式为 1 的整数矩阵群）：
\begin{equation*}
a_{I\mu}\rightarrow W_{IJ}a_{J\mu},\qquad W\in SL(n,Z) \tag{7.3.25}
\end{equation*}
由以上讨论我们看到，由 $(K_{1},\boldsymbol{q}_{1})$ 与 $(K_{2},\boldsymbol{q}_{2})$ 描述的两个 FQH 态是等价的（即它们属于同一个普适类（universality class）），如果存在 $W\in SL(n,Z)$ 使得
\begin{equation*}
\boldsymbol{q}_{2}=W\boldsymbol{q}_{1},\qquad K_{2}=WK_{1}W^{\top} \tag{7.3.26}
\end{equation*}
这是因为在变换 (7.3.25) 之下，由 $(K_{1},\boldsymbol{q}_{1})$ 描述的有效理论简单地变成由 $(K_{2},\boldsymbol{q}_{2})$ 描述的另一个有效理论。

我们想指出，在上述讨论中我们忽略了另一个拓扑量子数——自旋矢量（spin vector）。因此，式 (7.3.26) 的等价条件不适用于旋转不变系统。但是，式 (7.3.26) 确实适用于无序 FQH 系统，因为在无序系统中角动量不守恒，自旋矢量没有良好定义。关于自旋矢量的讨论可参见 Wen and Zee (1992c,d)。表 7.1 列出了一些常见的单层自旋极化 FQH 态的 $K$ 矩阵、电荷矢量 $\boldsymbol{q}$ 与自旋矢量 $\boldsymbol{s}$。

\subsection*{问题 7.3.3}
证明式 (7.3.24) 之下列出的四条性质。

\subsection{简单多层的分数量子霍尔态的有效理论}

\begin{keypoints}
\item $K$ 矩阵与多层 FQH 波函数。
\end{keypoints}

\begin{center}
\small 表 7.1\quad 某些简单的单层自旋极化 FQH 态的 $(K,\boldsymbol{q},\boldsymbol{s})$ 表达式\\[6pt]
\begin{tabular}{cccc}
\toprule
$\nu$ & $\boldsymbol{q}$ & $K$ & $\boldsymbol{s}$ \\
\midrule
$1/m$ & $(1)$ & $(m)$ & $\begin{pmatrix} m/2 \end{pmatrix}$ \\[10pt]
$1-1/m$ & $\begin{pmatrix}1\\0\end{pmatrix}$ &
$\begin{pmatrix}1 & 1\\ 1 & -(m-1)\end{pmatrix}$ &
$\begin{pmatrix}1/2\\ (1-m)/2\end{pmatrix}$ \\[18pt]
$2/5$ & $\begin{pmatrix}1\\0\end{pmatrix}$ &
$\begin{pmatrix}3 & -1\\ -1 & 2\end{pmatrix}$ &
$\begin{pmatrix}1/2\\ 1\end{pmatrix}$ \\[18pt]
$3/7$ & $\begin{pmatrix}1\\0\\0\end{pmatrix}$ &
$\begin{pmatrix}3 & -1 & 0\\ -1 & 2 & -1\\ 0 & -1 & 2\end{pmatrix}$ &
$\begin{pmatrix}1/2\\ 1\\ 1\end{pmatrix}$ \\
\bottomrule
\end{tabular}
\end{center}

到目前为止，我们只考虑了所谓的单层 FQH 态，其中二维电子气位于两种不同半导体的界面上。我们也可以制作具有多层界面的样品。置于强磁场之下时，不同层上的二维电子气可以形成多层 FQH 态。

用于构造分级态有效理论的同一方法，也可用来构造多层 FQH 态的有效理论。对多层 FQH 态而言，FQH 波函数与 $K$ 矩阵之间的联系变得非常透明。本节将集中讨论双层 FQH 态，不过推广到 $n$ 层 FQH 态是直截了当的。

我们想为下面这个简单的双层 FQH 态构造有效理论：
\begin{equation*}
\prod_{i<j}(z_{1i}-z_{1j})^{l}\prod_{i<j}(z_{2i}-z_{2j})^{m}\prod_{i,j}(z_{1i}-z_{2j})^{n}\,
\mathrm{e}^{-\frac{1}{4l_{B}^{2}}\left(\sum_{i}|z_{1i}|^{2}+\sum_{j}|z_{2j}|^{2}\right)}, \tag{7.3.27}
\end{equation*}
其中 $z_{Ii}$ 是第 $I$ 层中第 $i$ 个电子的复坐标。这里 $l$ 和 $m$ 是奇整数，使波函数与电子的费米统计相一致；$n$ 则可以是任意非负整数。上述波函数由 Halperin 首先提出，作为 Laughlin 波函数的推广 (Halperin, 1983)。这些波函数似乎能够解释双层 FQH 系统中观测到的部分主要 FQH 填充因子。

我们从第一层的单层 FQH 态出发，即 $\prod_{i<j}(z_{1i}-z_{1j})^{l}\mathrm{e}^{-\sum_{i}|z_{1i}|^{2}/4}$，它是一个 $1/l$ Laughlin 态，由如下有效理论描述
\begin{equation*}
\mathcal{L}=\left[-l\frac{1}{4\pi}a_{1\mu}\partial_{\nu}a_{1\lambda}\varepsilon^{\mu\nu\lambda}
+\frac{e}{2\pi}A_{\mu}\partial_{\nu}a_{\lambda}\varepsilon^{\mu\nu\lambda}\right] \tag{7.3.28}
\end{equation*}
其中 $a_{1\mu}$ 是描述第一层中电子密度和流的规范场。

考察式 (7.3.27) 的波函数可以看到，第二层中的一个电子束缚于第一层中的一个准空穴激发。这样的准空穴激发由 $n$ 个基本准空穴激发组成，并携带 $a_{1\mu}$ 荷 $-n$。一团准空穴由下面的有效理论描述：
\begin{equation*}
\mathcal{L}=-na_{1\mu}j^{\mu}+\text{动能/势能} \tag{7.3.29}
\end{equation*}
其中 $j^{\mu}$ 具有式 (7.3.2) 给出的形式。如前所述，在平均场理论中，如果忽略 $a_{1\mu}$ 场的响应，则式 (7.3.29) 简单地描述"磁场" $nb_{1}$ 中的一团玻色子，这里 $b_{1}=-\varepsilon_{ij}\partial_{i}a_{1j}$。现在我们想给式 (7.3.29) 中的每个准空穴附着一个（第二层中的）电子。这一操作有以下两个效果：由于电子携带电荷 $-e$，准空穴与电子的束缚态可以直接与电磁场 $A_{\mu}$ 耦合；并且束缚态表现得像一个费米子。束缚态的有效理论具有如下形式：
\begin{equation*}
\mathcal{L}=(eA_{\mu}-na_{1\mu})j^{\mu}+\text{动能/势能} \tag{7.3.30}
\end{equation*}
它现在（在平均场理论中）描述一团费米子。这些费米子感受到一个有效磁场 $-eB+nb_{1}$。

当第二层中的电子（即式 (7.3.30) 中的费米子）具有密度 $\frac{1}{m}(-eB+nb_{1})/2\pi$ 时（即它们具有有效填充因子 $1/m$），它们可以形成一个 $1/m$ Laughlin 态，对应于波函数中的 $\prod_{i<j}(z_{2i}-z_{2j})^{m}$ 部分。（注意这里 $eB<0$。）引入新规范场 $j^{\mu}=\frac{1}{2\pi}\partial_{\nu}a_{2\lambda}\varepsilon^{\mu\nu\lambda}$ 来描述式 (7.3.30) 中的费米子流 $j^{\mu}$，第二层中 $1/m$ 态的有效理论便具有形式
\begin{equation*}
\mathcal{L}=-m\frac{1}{4\pi}a_{2\mu}\partial_{\nu}a_{2\lambda}\varepsilon^{\mu\nu\lambda} \tag{7.3.31}
\end{equation*}
把式 (7.3.28)、式 (7.3.30) 与式 (7.3.31) 合在一起，我们就得到双层态的总有效理论，它具有式 (7.3.15) 给出的形式，其中 $K$ 矩阵与电荷矢量 $\boldsymbol{q}$ 由下式给出
\begin{equation*}
K=\begin{pmatrix} l & n\\ n & m \end{pmatrix},\qquad
\boldsymbol{q}=\begin{pmatrix}1\\1\end{pmatrix} \tag{7.3.32}
\end{equation*}
我们看到，$K$ 矩阵的矩阵元就是波函数中的幂次。该 FQH 态的填充因子仍由式 (7.3.19) 给出。

双层态中有两种（基本）准空穴激发。第一种通过在基态波函数上乘以 $\prod_{i}(\xi-z_{1i})$ 产生，第二种则由 $\prod(\xi-z_{2i})$ 产生。作为第一层与第二层这两个电子凝聚体中的涡旋，第一种准空穴由源项 $-a_{1\mu}j^{\mu}$ 产生，第二种准空穴由 $-a_{2\mu}j^{\mu}$ 产生。因此，双层态中一个一般的准粒子是若干第一类与第二类准空穴的束缚态，由式 (7.3.16) 描述。这些准粒子的量子数仍由式 (7.3.23) 给出。

\begin{center}
\small 表 7.2\quad 某些双层 FQH 态的 $(K,\boldsymbol{q},\boldsymbol{s})$ 表达式\\[6pt]
\begin{tabular}{cccc}
\toprule
$\nu$ & $\boldsymbol{q}$ & $K$ & $\boldsymbol{s}$ \\
\midrule
$1/m$ & $\begin{pmatrix}1\\1\end{pmatrix}$ &
$\begin{pmatrix}m & m\\ m & m\end{pmatrix}$ &
$\begin{pmatrix}m/2\\ m/2\end{pmatrix}$ \\[18pt]
$2/5$ & $\begin{pmatrix}1\\1\end{pmatrix}$ &
$\begin{pmatrix}3 & 2\\ 2 & 3\end{pmatrix}$ &
$\begin{pmatrix}3/2\\ 3/2\end{pmatrix}$ \\[18pt]
$1/2$ & $\begin{pmatrix}1\\1\end{pmatrix}$ &
$\begin{pmatrix}3 & 1\\ 1 & 3\end{pmatrix}$ &
$\begin{pmatrix}3/2\\ 3/2\end{pmatrix}$ \\[18pt]
$2/3$ & $\begin{pmatrix}1\\1\end{pmatrix}$ &
$\begin{pmatrix}3 & 0\\ 0 & 3\end{pmatrix}$ &
$\begin{pmatrix}3/2\\ 3/2\end{pmatrix}$ \\[18pt]
$2/3$ & $\begin{pmatrix}1\\1\end{pmatrix}$ &
$\begin{pmatrix}1 & 2\\ 2 & 1\end{pmatrix}$ &
$\begin{pmatrix}1/2\\ 1/2\end{pmatrix}$ \\
\bottomrule
\end{tabular}
\end{center}

一般地，(7.3.27) 类型的多层 FQH 态由一个矩阵元为整数、对角元为奇整数的 $K$ 矩阵描述。电荷矢量具有形式 $\boldsymbol{q}^{\top}=(1,1,\ldots,1)$。人们通常把双层 FQH 态 (7.3.27) 标记为 $(l,m,n)$。表 7.2 列出了一些简单双层态的 $K$ 矩阵、电荷矢量 $\boldsymbol{q}$ 与自旋矢量 $\boldsymbol{s}$。

由上述我们看到，(332) 双层态具有填充因子 2/5——这个填充因子也出现在单层分级态中。于是产生如下问题：双层 2/5 态与单层 2/5 态是否属于同一普适类？这个问题具有实验上的后果。我们可以设想这样的实验：从一个层间隧穿非常弱的系统中的 (332) 双层态出发，在保持填充因子不变的同时，把层间隧穿变得越来越强，双层态最终将变成单层 2/5 态。问题是：双层 (332) 态与单层 2/5 态之间的转变是平滑的渡越（crossover）还是相变？如果忽略自旋矢量，我们看到两个 2/5 态的 $K$ 矩阵和电荷矢量是等价的，因为它们通过一个 $SL(2,Z)$ 变换相联系。因此，在没有旋转对称性的情形下（此时自旋矢量没有良好定义），两个 2/5 态可以平滑地彼此转变。而当我们把自旋矢量包括进来时，两个 2/5 态并不等价，并且对旋转不变系统而言，它们被一个一级相变分隔开。

从表 7.2 我们还看到，存在两种不同的双层 2/3 态。当层内相互作用远强于层间相互作用时（真实样品中正是这种情形），基态波函数倾向于让同层电子之间具有更高阶的零点。因此，(330) 态的能量应当比 (112) 态低。单层 2/3 态的 $K$ 矩阵和电荷矢量与 (112) 态的等价，而与 (330) 态的不等价。因此，要把一个双层 2/3 态（即 (330) 态）变成单层 2/3 态，无论旋转对称性如何都必须经过一次相变。

双层 (mmm) 态是一个有趣的态，因为 $\det K=0$。Fertig (1989)、Brey (1990) 与 MacDonald 等 (1990) 研究了 (111) 态，发现了一个无能隙的集体模式。Wen 和 Zee (1992b) 研究了更一般的 (mmm) 态，并指出 (mmm) 态（在无层间隧穿时）自发破缺一个 $U(1)$ 对称性，是一种中性超流体。更详细的讨论（包括其实验含义）可参见 Wen and Zee (1993)、Murphy et al. (1994) 与 Yang et al. (1994)。

\subsection*{问题 7.3.4}
让我们首先忽略自旋矢量 $\boldsymbol{s}$。证明表 7.1 中单层 2/5 态的 $(K,\boldsymbol{q})$ 与表 7.2 中双层 2/5 态的 $(K,\boldsymbol{q})$ 等价。再证明，这两个态的 $(K,\boldsymbol{q},\boldsymbol{s})$ 却不等价。（在存在自旋矢量 $\boldsymbol{s}$ 时，等价关系变为 $\boldsymbol{s}_{2}=W\boldsymbol{s}_{1}$、$\boldsymbol{q}_{2}=W\boldsymbol{q}_{1}$ 与 $K_{2}=WK_{1}W^{+}$。）

\subsection*{问题 7.3.5}
考虑双层 (lmn) 态 (7.2.6)。利用有效理论 (7.3.15) 与式 (7.3.32) 求出填充因子。并求出第一层与第二层中准空穴的分数电荷与分数统计。（可以把结果与问题 7.2.5 的结果比较。）

\subsection*{问题 7.3.6}
假设双层 (mmm) 态的有效理论具有形式
\begin{equation*}
\mathcal{L}=-\frac{1}{4\pi}K_{IJ}a_{I\mu}\partial_{\nu}a_{J\lambda}\varepsilon^{\mu\nu\lambda}
+\frac{1}{2g_{1}}\pmb{e}_{I}\cdot\pmb{e}_{I}-\frac{1}{2g_{2}}b_{I}b_{I}
\end{equation*}
其中 $\pmb{e}_{I}$ 与 $b_{I}$ 分别是 $a_{I\mu}$ 的"电场"与"磁场"。
\begin{enumerate}
\item 证明 (mmm) 态正如超流体一样具有无能隙激发。求出无能隙激发的速度。
\item 同样像超流体一样，(mmm) 态包含能量为 $\gamma\ln L$ 的涡旋激发，这里 $L$ 是系统的线度。求出能量最小的涡旋的 $\gamma$ 值。（提示：在 (mmm) 态中，$l_{I}$ 的 $a_{I\mu}$ 荷仍量子化为整数。）
\item 如果我们确实把 (mmm) 态当作超流体，你能辨认出破缺的对称性和序参量吗？
\end{enumerate}

\section{分数量子霍尔液体中的边缘激发}

\begin{keypoints}
\item FQH 态\emph{总是}具有无能隙的边缘激发。
\item FQH 态的边缘激发构成手征 Luttinger 液体。
\item 边缘激发的结构由体拓扑序决定。
\end{keypoints}

由于电子之间排斥性的相互作用和强关联，尽管第一朗道能级只是部分填充，QH 液体仍是一个不可压缩态。QH 态中的所有体激发都具有有限的能隙。QH 态与绝缘体非常相似：两者都具有有限能隙和短程的电子传播子。正是由于这种相似性，人们曾对如下事实感到困惑：QH 系统的输运性质显然与普通绝缘体截然不同。Halperin 首先指出，IQH 态含有无能隙的边缘激发 (Halperin, 1982)。IQH 态的非平庸输运性质正来自这些无能隙边缘激发 (Halperin, 1982; Trugman, 1983; MacDonald and Streda, 1984; Streda et al., 1987; Buttiker, 1988; Jain and Kivelson, 1988a,b)。例如，对一个 IQH 样品做两端测量，只有当源与汇通过边缘连通时才可能给出有限电阻。如果源与汇不被任何边缘连通，那么两端测量在零温下将给出无穷大电阻——这是一个与绝缘体十分相似的结果。Halperin 还研究了 IQH 态边缘激发的动力学性质，发现边缘激发由一个手征的一维费米液体理论描述。

鉴于 FQH 态与 IQH 态输运性质的相似性，一个自然的猜想是：FQH 态的输运同样由边缘激发支配 (Beenakker, 1990; MacDonald, 1990)。然而，由于 FQH 态本质上是多体态，FQH 态中的边缘激发无法通过填充单粒子能级来构造。换言之，FQH 态的边缘激发不应当用费米液体来描述。因此，我们需要一种全新的方法来理解 FQH 液体边缘态的动力学性质 (Wen, 1992, 1995)。

在下一节中，我们将用费米液体理论研究 $\nu=1$ IQH 态的边缘激发。然后，我们将利用流代数（或玻色化，见 Tomonaga, 1950; Luttinger, 1963; Coleman, 1975）来建立 FQH 态边缘激发的理论。

\subsection{整数量子霍尔边缘态的费米液体理论}

\begin{keypoints}
\item IQH 液体的基态通过填充单粒子能级得到。因此，低能边缘激发可以通过以略有不同的方式填充这些单粒子能级来获得。
\end{keypoints}

考虑均匀磁场 $B$ 中的一团无相互作用电子气。在对称规范下，角动量是一个好量子数。能量与动量的共同本征态构成一个环状分布（见图 7.10）。因此，在光滑的圆形势 $V(r)$ 存在时，单粒子能级如图 7.15 所示。基态通过填充能量最低的 $N$ 个能级得到。这样的态对应于一个具有均匀密度的圆形液滴（见图 7.11）。

%FIG{7.15}: {光滑势 $V(r)$ 存在时头三个朗道能级中的能级。这里 $m$ 是能级的角动量。化学势 $\mu$ 以下的能级各被一个电子占据，从而给出 IQH 态。改变边缘附近的占据数，便产生 IQH 态的无能隙边缘激发。}

在零朗道能级中，角动量 $m$ 的态的能量为
\begin{equation*}
E_{m}=\frac{1}{2}\hbar\omega_{c}+V(r_{m})
\end{equation*}
其中 $r_{m}=\sqrt{2m}\,l_{B}$ 是 $m$ 态的半径。从图 7.15 我们看到，体激发具有有限的能隙 $\hbar\omega_{c}$。边缘激发可以这样产生：把一个电子从边缘附近的 $m$ 态移到 $m+1$ 态（见图 7.15）。由于在热力学极限 $r_{m}\to\infty$ 下 $r_{m+1}-r_{m}\to0$，边缘激发是无能隙的。

由于 $m$ 态具有半径 $r_{m}$，我们也可以把角动量 $m$ 看作沿边缘的动量 $k=m/r_{m}$。这使我们能够把 $E_{m}$ 当作一个能量--动量关系
\begin{equation*}
E(k)=\frac{1}{2}\hbar\omega_{c}+V(2l_{B}^{2}k)
\end{equation*}
用 $E(k)$ 来表示，我们的电子系统可以用一个无相互作用的哈密顿量描述：
\begin{equation*}
H=\sum_{k}E(k)c_{k}^{\dagger}c_{k} \tag{7.4.1}
\end{equation*}
上述哈密顿量描述一个一维手征费米液体（chiral Fermi liquid）。我们称它为手征费米液体，是因为它只有一个费米点，并且所有低能激发都朝同一个方向传播。\footnote{相比之下，通常的一维费米液体有两个费米点，同时包含左移与右移两种成分。}我们由此得出结论：$\nu=1$ IQH 态的边缘激发由一维手征费米液体 (7.4.1) 描述。

%FIG{7.16}: {(a) 能量--角动量关系 $E_{m}$ 可以看作能量--动量关系 $E(k)$。边缘激发可以看作一维手征费米液体的激发，其中只含右移成分。(b) 通常的一维费米液体的色散关系同时包含右移与左移成分。}

\subsection*{问题 7.4.1}
证明 IQH 边缘激发的速度 $v=\partial E(k)/\partial k$ 由 $cE/B$ 给出，其中 $c$ 是光速，$E=-\partial V(r)/\partial r$ 是约束势 $V(r)$ 在边缘处产生的电场。设 $\nu=1$ IQH 液滴包含 $N$ 个电子，求出手征费米液体的费米动量 $k_{F}$。

\subsection*{问题 7.4.2}
一个 $\nu=1$ IQH 态被圆形光滑势 $V(r)$ 约束，包含 $N$ 个电子。求出基态的总角动量 $M_{0}$。求出总角动量为 $M_{0}+m$ 的低能粒子--空穴激发的数目，其中 $m=-1$ 及 $m=1,2,3,4,5$。证明在热力学极限下，具有相同总角动量的激发具有相同的能量。

\subsection*{问题 7.4.3}
考虑一个由 $H=\sum_{k}(v_{1}k\,c_{1,k}^{\dagger}c_{1,k}+v_{2}k\,c_{2,k}^{\dagger}c_{2,k})$ 描述的一维无相互作用费米子系统。混合项 $\sum_{k}(\gamma c_{1,k}^{\dagger}c_{2,k}+\text{h.c.})$ 是一个相关微扰。证明：当系统同时具有右移与左移成分（即 $v_{1}v_{2}<0$）时，这个相关微扰会打开一个有限的能隙，并剧烈地改变系统的低能性质。然而，对于所有激发都朝同一方向运动的手征费米液体（即 $v_{1}v_{2}>0$），这个相关微扰不诱导任何能隙。事实上，我们认为任何微扰都无法在一维手征费米液体中诱导出能隙。一维手征费米液体中的无能隙激发对所有微扰都是稳健的，并且在拓扑上是稳定的。

\subsection{流体动力学方法——$1/m$ Laughlin 态}

\begin{keypoints}
\item 边缘激发来自流代数（Kac--Moody 代数）。
\item 无法通过填充单粒子能级来构造 FQH 边缘激发。
\end{keypoints}

由于 $1/m$ Laughlin FQH 态不能通过填充单粒子能级得到，我们不能用上一节的图像来构造 $1/m$ Laughlin 态的边缘激发。由于不知道如何从相互作用电子的哈密顿量导出 FQH 边缘理论，下面我们转而尝试猜测一个低能有效理论。

猜测边缘激发动力学理论的最简单方法是使用流体动力学方法。在这种方法中，人们利用这样一个事实：FQH 态是不可压缩、无旋的液体，不含低能体激发。因此，唯一的低能激发（在体能隙之下）是 FQH 液滴形变产生的表面波。这些表面波被认定为 FQH 态的边缘激发。

在流体动力学方法 (Wen, 1992) 中，我们首先研究 FQH 液滴上表面波的经典理论，然后把经典理论量子化，得到边缘激发的量子描述。令人惊叹的是，由经典理论得到的这个简单的量子描述，在低能下给出了边缘激发的完整描述，并使我们能够计算沿边缘的电子与准粒子传播子。

考虑一个填充因子为 $\nu$、被势阱约束的 FQH 液滴。由于电导非零，势阱的电场产生一个沿边缘流动的持续电流：
\begin{equation*}
\pmb{j}=\sigma_{xy}\hat{z}\times\pmb{E},\qquad \sigma_{xy}=\nu\frac{e^{2}}{h}
\end{equation*}
这意味着边缘附近的电子以速度
\begin{equation*}
v=\frac{E}{B}c
\end{equation*}
沿边缘漂移，这里 $c$ 是光速。于是，边缘波（边缘的形变）也以速度 $v$ 传播。让我们用一维密度 $\rho(x)=nh(x)$ 来描述边缘波，其中 $h(x)$ 是边缘的位移，$x$ 是沿边缘的坐标，而 $n=\nu/2\pi l_{B}^{2}$ 是体内的二维电子密度。（这里 $l_{B}=\sqrt{c/eB}$ 是磁长度。）我们看到，边缘波的传播由下面的波动方程描述：
\begin{equation*}
\partial_{t}\rho+v\partial_{x}\rho=0 \tag{7.4.2}
\end{equation*}
注意，边缘波总是朝一个方向传播；不存在朝相反方向传播的波。

边缘波的哈密顿量（即能量）由下式给出
\begin{equation*}
H=\int\mathrm{d}x\,\frac{1}{2}e\rho E h=\int\mathrm{d}x\,\pi\frac{v}{\nu}\rho^{2} \tag{7.4.3}
\end{equation*}
其中我们用到了 $h=\rho/n=\rho\frac{2\pi c}{eB}$。在动量空间中，式 (7.4.2) 与式 (7.4.3) 可以改写为
\begin{equation*}
\dot{\rho}_{k}=-\mathrm{i}vk\rho_{k},\qquad
H=2\pi\frac{v}{\nu}\sum_{k>0}\rho_{-k}\rho_{k} \tag{7.4.4}
\end{equation*}
其中 $\rho_{k}=\int\mathrm{d}x\,\frac{1}{\sqrt{L}}\mathrm{e}^{\mathrm{i}kx}\rho(x)$，而 $L$ 是边缘的长度。我们发现，如果把 $\rho_{k}|_{k>0}$ 认作"坐标"，并把 $\pi_{k}=\mathrm{i}2\pi\rho_{-k}/\nu k$ 认作相应的正则"动量"，那么标准哈密顿方程
\begin{equation*}
\dot{q}=\frac{\partial H}{\partial p},\qquad \dot{p}=-\frac{\partial H}{\partial q}
\end{equation*}
将重现运动方程 $\dot{\rho}_{k}=\mathrm{i}vk\rho_{k}$。这使我们得以认定正则"坐标"与"动量"。有趣的是，位移 $h(x)$ 同时包含"坐标"与"动量"。这是边缘波的手征性质所致。

知道了正则坐标与动量之后，量子化经典理论就很容易了。我们只需把 $\rho_{k}$ 和 $\pi_{k}$ 看作满足 $[\rho_{k},\pi_{k'}]=\mathrm{i}\delta_{kk'}$ 的算符。于是，量子化之后我们得到
\begin{align*}
[\rho_{k},\rho_{k'}]&=\frac{\nu}{2\pi}k\delta_{k+k'},\qquad
k,k'=\text{整数}\times\frac{2\pi}{L} \tag{7.4.5}\\
H&=2\pi\frac{v}{\nu}\sum_{k>0}\rho_{-k}\rho_{k}
\end{align*}
上述代数称为 (U(1)) Kac--Moody（K--M）代数 (Kac, 1983; Goddard and Olive, 1985, 1986)。类似的代数也曾出现在 Tomonaga 模型中 (Tomonaga, 1950; Luttinger, 1963; Coleman, 1975)。注意，式 (7.4.5) 简单地描述了一组退耦谐振子（由 $(\rho_{k},\rho_{-k})$ 生成）。因此，式 (7.4.5) 是一个一维自由声子理论（只有单一分支）。对于 $k>0$，$\rho_{k}^{\dagger}=\rho_{-k}$ 产生一个动量为 $k$、能量为 $vk$ 的声子，而 $\rho_{k}$ 湮灭一个声子。式 (7.4.5) 为 Laughlin 态的低能边缘激发提供了完整的描述。

这里考虑的边缘激发不改变系统的总电荷，因而是中性的。下面我们将讨论带电激发，并由 K--M 代数 (7.4.5) 计算电子传播子。

低能电荷激发显然对应于把电子加到边缘上（从边缘移走）。这些带电激发携带整数电荷，并且
